\section{Finite Morita reduction for \texorpdfstring{$\RR{8}{1660}$}{R6\_8\_1660}}\label{sec:1660-morita-reduction}
\begin{proposition}[A general formal-codegree test for modules]
Let $\mathcal F$ be pseudo-unitary with commutative Grothendieck ring.
For an indecomposable module category, let $m_\chi$ be the multiplicity
of the character $\chi$ in its complexified based module, and let
$f_\chi$ be the formal codegree. Then
\[
 \sum_\chi\frac{m_\chi}{f_\chi}\leq1.
\]
If equality holds, the Morita dual has rank $\sum_\chi m_\chi^2$,
and its action is faithful and equals the full commutant. In this
case its simple action matrices must give a square nonnegative
integral Gram factorization of the linking-category matrix below.
Failure of that bounded factorization is an obstruction to the based
module being a module category.
\end{proposition}
\begin{proof}
Ostrik's module-character reciprocity identifies $m_\chi$ with the
multiplicity in the module's Lagrangian algebra of the central simple
of dimension $\dim(\mathcal F)/f_\chi$. Positivity and the dimension
of that algebra give the inequality. Equality exhausts its support
inside the canonical algebra. Its endomorphism algebra is consequently
$\bigoplus_\chi M_{m_\chi}(\mathbb C)$ and acts faithfully on the
module Hom space. The next proof details these identifications and
their precise primary source. The linking associativity identity then
gives the stated Gram factorization.
\end{proof}

This note gives necessary module-category conditions, not an existence or
nonexistence theorem. No braiding on $\mathcal C$ or its relative center is assumed.
Let $\mathcal C$ have simples $1,g,X,gX,\sigma,T=\sigma X$ and rules
\[
g^2=1,\quad g\sigma=\sigma,\quad\sigma^2=1+g,\qquad
X^2=1+\bigl(3(1+g)+4\sigma\bigr)X.
\]
Put $d=3+2\sqrt2+2\sqrt3+\sqrt6$, so $\FPdim\mathcal C=4(1+d^2)$.
All matrices below use the action on the basis of simple objects of a
finite semisimple module category.

\begin{lemma}
Every indecomposable based module over the Ising fusion ring, whose
simple action matrices obey duality transpose, is its three-object regular
module. Consequently every $\mathcal C$-module category restricts to a direct sum
of regular Ising module categories at the level of based modules.
\end{lemma}
\begin{proof}
Write $H=N_\sigma$ and $G=N_g$. Then $G$ is a permutation involution,
$H=H^{\mathsf T}\geq0$, $GH=H$, and $H^2=I+G$. An indecomposable
component of $H$ has spectral radius $\sqrt2$: this follows either from
its dimension vector or directly from these identities and $GH=H$.
There are no entries of $H$ greater than one. There are no loops: a loop
connected to any different vertex has norm at least the golden ratio,
while an isolated loop contradicts $H^2=I+G$. No vertex has degree three
or more, and a path with four vertices already has norm greater than
$\sqrt2$. Thus every connected component is the three-vertex path
$A_3$; the possibilities of one or two vertices contradict $H^2=I+G$.
Its two leaves are exchanged by $G$ and its middle vertex is fixed.
These are the regular Ising matrices.
\end{proof}

\begin{theorem}[Uniform module rank]
Every indecomposable finite semisimple $\mathcal C$-module category has exactly
six simple objects. Its restriction to $\mathcal I=\langle g,\sigma\rangle$
consists of two regular Ising blocks.
\end{theorem}
\begin{proof}
Initially suppose there are $n$ blocks. Order each block as $1,g,\sigma$.
The commutant of the two regular Ising matrices consists exactly of
\[
J(a,b,c)=\begin{pmatrix}a&b&c\\b&a&c\\c&c&a+b\end{pmatrix}.
\]
Since $X$ commutes with $g,\sigma$, and is self-dual, its action has form
$N_X=P\otimes I+Q\otimes N_g+R\otimes N_\sigma$,
where $P,Q,R$ are symmetric nonnegative integral $n$ by $n$ matrices.
Set $U=P+Q$, $V=P-Q$. The $X^2$ relation gives
\begin{equation}
V^2=I,\quad U^2+2R^2=I+6U+8R,\quad UR+RU=6R+4U. \label{eq:1660-morita-relations}
\end{equation}
The matrix $M=U+\sqrt2R$ therefore satisfies
\begin{equation}
M^2-(6+4\sqrt2)M-I=0. \label{eq:1660-morita-quadratic}
\end{equation}
It is irreducible: otherwise a union of Ising blocks would be invariant
under all the generators of $K_0(\mathcal C)$, contradicting indecomposability.
Hence its positive Perron eigenvalue $d$ is simple. The polynomial in
\eqref{eq:1660-morita-quadratic} is irreducible over $\mathbb Q(\sqrt2)$: its discriminant is
$72+48\sqrt2=(4\sqrt3+2\sqrt6)^2$, whose square root does not belong
to $\mathbb Q(\sqrt2)$. Thus the characteristic polynomial of $M$ is a
power of \eqref{eq:1660-morita-quadratic}; simplicity of the Perron root forces that power to be one.
Therefore $n=2$.
\end{proof}

\begin{theorem}[Complete based-module list]
Up to interchanging the two Ising blocks, the possible matrices $U,R$
are exactly
\[
U=\begin{pmatrix}a&b\\b&6-a\end{pmatrix},\qquad
R=\begin{pmatrix}c&f\\f&4-c\end{pmatrix},
\]
with the following four rows. These are based-module possibilities;
their associativity constraints have not been solved.
\begin{center}
\begin{tabular}{c|rrrr|l}
type&$a$&$b$&$c$&$f$&$V=P-Q$\\\hline
I&0&1&0&0&signed transposition\\
II&2&3&2&2&signed transposition\\
III&1&2&0&1&$\operatorname{diag}(\epsilon_1,\epsilon_2)$\\
IV&1&2&1&2&$\operatorname{diag}(\epsilon_1,\epsilon_2)$
\end{tabular}
\end{center}
Here $\epsilon_i\in\{1,-1\}$ and $P=(U+V)/2$, $Q=(U-V)/2$.
The two signed transpositions give equivalent modules by swapping the
two leaves in one Ising block. Thus there are ten labelled-$X$ based
modules, or six after additionally allowing the ring automorphism
$X\leftrightarrow gX$.
\end{theorem}
\begin{proof}
Equation \eqref{eq:1660-morita-quadratic} and $n=2$ give $\operatorname{tr}U=6$ and
$\operatorname{tr}R=4$. The matrix $V$ is symmetric integral and
$V^2=I$, so it is a signed permutation. It must satisfy
$V\equiv U\pmod2$, with $U\pm V\geq0$ entrywise.
Put $w=a-3$, $s=c-2$. The two equations \eqref{eq:1660-morita-relations} become
\begin{equation}
ws+bf=6,\qquad w^2+b^2+2(s^2+f^2)=18. \label{eq:1660-morita-small-squares}
\end{equation}
All variables are integral; $|w|\leq3$, $|s|\leq2$, and $b,f\geq0$.
Cauchy--Schwarz implies $3\leq s^2+f^2\leq6$.
The only sums of two integral squares in this interval are $4$ and $5$.
For value $4$, the squared norm of $(w,b)$ is $10$; for value $5$, it
is $8$. Substitution in $ws+bf=6$ gives exactly the four rows and their
block exchanges. Parity gives precisely the displayed options for $V$.
Conversely these matrices are nonnegative integral and satisfy every
fusion relation and transpose relation. Their $X$ action connects the
two Ising blocks, so they are indecomposable based modules.
\end{proof}

\begin{corollary}[Four small algebra supports]
Every indecomposable $\mathcal C$-module category has a connected separable
internal-end algebra on one of the following underlying objects:
\begin{equation}
1,\qquad 1+X+gX+2T,\qquad 1+X\ \text{or }1+gX,
\qquad 1+X+T\ \text{or }1+gX+T. \label{eq:1660-morita-supports}
\end{equation}
Consequently these supports exhaust the possible Morita classes, with
no assertion that the nontrivial supports admit algebra structures or
that different structures give different Morita classes.
\end{corollary}
\begin{proof}
Choose a leaf $m$ in the first Ising block of each row. The coefficient
of a simple $Z$ in $\underline{\mathrm{End}}(m)$ is $(N_Z)_{mm}$.
Thus this internal end is
$1+P_{11}X+Q_{11}gX+R_{11}T$. Substitution gives \eqref{eq:1660-morita-supports}.
The standard internal-Hom construction supplies a connected separable
algebra and realizes the original module category as its modules.
The first row has internal end $1$ and is necessarily the regular
module category, even categorically.
\end{proof}

The Perron dimension ratio of the second Ising block to the first is,
respectively,
\[
d,\quad 3-2\sqrt2+2\sqrt3-\sqrt6,\quad
\sqrt2+\sqrt3,\quad(\sqrt2+\sqrt6)/2.
\]
The dimension quotients $\FPdim\mathcal C/\FPdim Q$ for the three nontrivial
supports in \eqref{eq:1660-morita-supports} are, respectively,
\[
144-96\sqrt2+80\sqrt3-56\sqrt6,
\qquad24+8\sqrt6,\qquad12+4\sqrt3.
\]
All are algebraic integers, so elementary dimension divisibility does
not exclude these targets. The two choices of the second diagonal sign
in types III and IV are not removed by the chosen algebra support.
They may encode different algebra/module associativity structures.

\begin{theorem}[Full-center saturation and ranks of all Morita duals]
Use the established reduction to the pseudo-unitary form of $\mathcal C$.
For any indecomposable module category $\mathcal M$, the dual fusion category
$\mathcal D=\operatorname{Fun}_{\mathcal C}(\mathcal M,\mathcal M)$ has rank six or eight.
More precisely,
\[
\begin{array}{c|c|c}
\text{matrix }V&\operatorname{rank}\mathcal D&
K_0(\mathcal D)\otimes\mathbb C\\\hline
\operatorname{spec}(V)=\{1,-1\}&6&\mathbb C^6\\
V=I\text{ or }-I&8&\mathbb C^4\oplus M_2(\mathbb C).
\end{array}
\]
The action of $K_0(\mathcal D)\otimes\mathbb C$ on
$K_0(\mathcal M)\otimes\mathbb C$ is faithful and is the full commutant of
the $K_0(\mathcal C)$ action.
\end{theorem}
\begin{proof}
Write $D=\FPdim\mathcal C$ and let $A_\chi$ denote the six simple summands
of the canonical algebra $I_{\mathcal C}(1)$. They occur once, because
$K_0(\mathcal C)$ is commutative. Ostrik's Theorem 2.13 gives
$\dim A_\chi=D/f_\chi$.
Let $L_\mathcal M$ be the Lagrangian algebra in $Z(\mathcal C)$ corresponding to
$\mathcal M$; under the Morita equivalence of centers it is $I_{\mathcal D}(1)$.
The precise module-character reciprocity needed here is given by
Ostrik's Remark 2.19: the bimodule $K_0(\mathcal M)\otimes\mathbb C$ is
identified with
\[
\operatorname{Hom}_{Z(\mathcal C)}(I_{\mathcal C}(1),L_\mathcal M)
\]
over the corner algebras $\operatorname{End}(I_{\mathcal C}(1))$ and
$\operatorname{End}(L_\mathcal M)$. This follows by applying that remark to
the linking multi-fusion category of $\mathcal C$ and $\mathcal M$.
Its twist hypothesis is automatic in the pseudo-unitary case, as stated
at the end of the same remark. Thus the multiplicity of $\chi$ in the
module character is exactly $[L_\mathcal M:A_\chi]$.

The four characters with $g=1$ each occur once: each of the two Ising
spin evaluations gives a two-dimensional matrix with its two distinct
roots. The two characters with $g=-1$, $\sigma=0$, $X=\pm1$ occur
with multiplicities $m_+,m_-$ equal to the eigenvalue multiplicities of
$V$; hence $m_++m_-=2$. Both have formal codegree $4$. Consequently
the dimensions of all the already specified constituents of $L_\mathcal M$ sum
to exactly $D$, just as those of $I_{\mathcal C}(1)$ do. Since
$\dim L_\mathcal M=D$ and all dimensions are positive, there are no additional
constituents. Therefore
\[
L_\mathcal M=\bigoplus_{\chi:g=1}A_\chi\ \oplus\ m_+A_+\oplus m_-A_-,
\qquad
\dim\operatorname{End}(L_\mathcal M)=4+m_+^2+m_-^2.
\]
Ostrik's Corollary 2.16 identifies this endomorphism algebra with the
complexified Grothendieck algebra of $\mathcal D$, proving the table.
Every simple constituent of $L_\mathcal M$ also occurs in $I_{\mathcal C}(1)$.
Consequently the action of $\operatorname{End}(L_\mathcal M)$ on the displayed
Hom space sees every matrix block and is faithful. Its block
decomposition also shows that it is the full commutant.
\end{proof}

\paragraph{A first forced Morita-dual ring.}
For type III with $V=\operatorname{diag}(1,-1)$, the rank-six dual,
if the module category exists, has the original fusion ring
$\RR{8}{1660}$. Here is a short exact reconstruction.
Write $N_x$ for the six known action matrices on $\mathcal M$, and $Q_y$ for
the unknown right action matrices of the dual. Associativity in the
linking multi-fusion category and Frobenius reciprocity imply
\begin{equation}
\sum_x N_x[a,b]N_x[c,d]
 =\sum_y Q_y[c,a]Q_y[d,b]. \label{eq:1660-morita-linking-gram}
\end{equation}
The central fermion of $\mathcal C$ gives a dual invertible whose action is
$N_g$, so the two matrices $I,N_g$ are fixed. Choose a leaf $a$ in the
first Ising block. The diagonal entry of \eqref{eq:1660-morita-linking-gram} at the coordinate $(a,a)$
is $2$; after removing $I$, precisely one further simple matrix $Q_Y$
has entry $1$ there. It is therefore the corresponding Gram column
minus the column of $I\otimes I$. Its $g$-translate is another simple.
After removing these four matrices, the residual Gram matrix has rank
two. At coordinate $(a,\sigma a)$ its diagonal is $2$, so both remaining
matrices have entry $1$. That Gram column is their sum $v$; if $H$ is
the residual Gram matrix, $2H-vv^{\mathsf T}$ is the outer product of
their difference. This determines them up to exchange.
The resulting matrices, called $Q_Z,Q_s$, satisfy
\[
s^2=1+g,\qquad sY=Z,\qquad
Y^2=1+3Y+3gY+4Z.
\]
The exact script \texttt{reconstruct\_type\_iii.py} verifies \eqref{eq:1660-morita-linking-gram},
nonnegative integrality, linear independence, and all fusion products.
Thus this branch is a possible Morita self-partner, not an exclusion.

\begin{theorem}[Complete integral Morita Gram reduction]
The two type-IV based modules with $V=I$ or $V=-I$ cannot be module
categories. For the other types, equation \eqref{eq:1660-morita-linking-gram}, the central fermion,
and the based-ring axioms leave exactly the following action bases.
Set
\[
a_0=2+\sqrt2+\sqrt3+\sqrt6,\quad
b_0=2+2\sqrt2+\sqrt3+\sqrt6,\quad
c_0=5+4\sqrt2+3\sqrt3+2\sqrt6.
\]
\begin{center}\begin{tabular}{c|c|l}
module type&number&dimensions of forced dual rings\\\hline
II&1&$1,1,\sqrt2,d,d,\sqrt2d$\\
III, mixed signs&1&$1,1,\sqrt2,d,d,\sqrt2d$\\
III, equal signs&2&$1,1,d,d,a_0,a_0,b_0,b_0$\\
IV, mixed signs&2&$1,1,\sqrt2,d,d,\sqrt2d$ or $1,1,1+\sqrt3,a_0,a_0,c_0$\\
IV, equal signs&0&none.
\end{tabular}\end{center}
The three occurrences of the first dimension list have the original
ring. The second rank-six ring is self-dual with commuting generators
$g,S,Y$ satisfying
\[
g^2=1,\quad gS=S,\quad S^2=1+g+2S,\qquad
Y^2=1+(1+g+2S)Y.
\]
Its maximal structure constant is $16$.
The two rank-eight rings have maximal structure constant $3$ and are
not isomorphic: in one the $b_0$ pair is a dual pair and all other simples
are self-dual; in the other the $a_0$ pair is a dual pair and all other
simples are self-dual.
\end{theorem}
\begin{proof}
Here is a complete finite calculation with explicit bounds. It is
implemented in \texttt{reconstruct\_dual\_actions.py}; its output is
\texttt{dual\_actions\_receipt.json}. Flatten action matrices to vectors.
Equation \eqref{eq:1660-morita-linking-gram} gives a known nonnegative symmetric $36$ by $36$ matrix
$K=\sum_x N_x\otimes N_x$. Subtract the outer products of the known
matrices $I,N_g$. For type III also subtract the two matrices forced in
the preceding reconstruction. The residual matrix $H$ has rank $r$,
equal to the number of remaining dual simples.

Choose $r$ coordinates whose principal submatrix $B$ is invertible.
Any remaining simple column $q$ is determined by its coordinates $v$
there: $q=H_{*,J}B^{-1}v$. Because there are exactly $r$ remaining
columns, their square factorization of $B$ implies
\[
v^{\mathsf T}B^{-1}v=1,
\qquad 0\leq v_i\leq\lfloor\sqrt{B_{ii}}\rfloor.
\]
Conversely, an orthogonal basis of $r$ such vectors, lifting to
nonnegative integral columns, gives the whole Gram factorization.
This is a bounded integer ellipsoid enumeration, followed by a clique
search for orthogonal bases. The numbers of candidate columns and
orthogonal bases for rows II, III mixed, III equal, IV mixed, IV equal
are respectively
\[
(4,1),\quad(2,1),\quad(10,2),\quad(12,3),\quad(12,0).
\]
Testing nonnegative integral products, closure under transpose, and the
unit coefficient duality rule leaves respectively $1,1,2,2,0$ fusion
action bases. Faithfulness proved above makes their multiplication the
actual dual fusion ring. The receipt records every resulting matrix
and every structure constant, so the counts and identifications require
no numerical optimizer. The signs $V\mapsto -V$ amount to relabelling
$X$ and $gX$, which leaves $K=\sum_xN_x\otimes N_x$ unchanged.
\end{proof}

\begin{corollary}[A Morita second-indicator exclusion]
The rank-eight ring in which the $a_0$ pair is nonself-dual cannot occur
as a Morita dual of $\mathcal C$. In the other rank-eight possibility every
self-dual simple has second indicator $+1$.
\end{corollary}
\begin{proof}
For a pseudo-unitary category $\mathcal F$, put
$W_2(\mathcal F)=\sum_Z\nu_2(Z)\dim Z$.
The center indicator formula gives
$\tau_2(Z(\mathcal F))=\dim(\mathcal F)W_2(\mathcal F)$.
Positive spherical structures and the twist are preserved by Morita
equivalence of centers, so $W_2$ is a Morita invariant here. The
established indicators of $\mathcal C$ give
\[
W_2(\mathcal C)=2-\sqrt2+(2+\sqrt2)d
       =12+6\sqrt2+6\sqrt3+4\sqrt6.
\]
If the $a_0$ pair is nonself-dual, the possible indicator sum has form
$u+s_d d+s_b b_0$, where $u\in\{0,2\}$ and
$s_d,s_b\in\{-2,0,2\}$. Comparing $\sqrt3$ and $\sqrt6$
coefficients forces $s_d=s_b=2$. The resulting $\sqrt2$ coefficient
is $8$, not $6$, a contradiction. This argument does not need a prior
choice of the indicator of the nontrivial invertible.
If instead the $b_0$ pair is nonself-dual, the same coefficient
comparison gives $s_d=s_a=2$ and then $u=2$.
Indeed $W_2(\mathcal C)=2+2d+2a_0$.
\end{proof}

\begin{corollary}[No ordinary generalized-near-group Morita model]
No category with this ring can be Morita equivalent to a generalized
near-group category in which invertibles act transitively on all
noninvertible simples. The same holds after any finite chain of finite
group extensions, equivariantizations and inverse such operations
starting from an ordinary generalized-near-group category.
\end{corollary}
\begin{proof}
For a generalized-near-group category the noninvertibles have common
dimension $z$, and $z^2=k+mz$ for integers $k>0,m\geq0$. Its global
dimension lies in $\mathbb Q(z)$, of degree at most two. Morita
equivalence preserves FP dimension; extensions and equivariantizations
multiply it by positive integers, and their inverses divide by integers.
In contrast $\FPdim\mathcal C=144+96\sqrt2+80\sqrt3+56\sqrt6$ has degree four:
none of the three nontrivial automorphisms of
$\mathbb Q(\sqrt2,\sqrt3)$ fixes it.
\end{proof}

\paragraph{Remaining gap.}
The displayed nimreps do not supply their module associators.
The full-center and Gram arguments now leave only three possible
Morita-dual based rings: the original rank-six ring, the displayed
rank-six $\operatorname{Ad}(E_6)$ extension, and one rank-eight ring.
Their existence, and the associators of the candidate modules, remain
unproved. The present note excludes module/dual alternatives but does
not exclude or construct $\RR{8}{1660}$.

\paragraph{Background sources.}
The internal-Hom/algebra description of semisimple module categories is
standard; see V. Ostrik, \emph{Module categories, weak Hopf algebras and
modular invariants}, \href{https://arxiv.org/abs/math/0111139}{arXiv:math/0111139}.
For module traces in the pseudo-unitary case and their symmetric
Frobenius algebras, see G. Schaumann, \emph{Traces on module categories
over fusion categories}, especially Section 5 and Section 6,
\href{https://arxiv.org/abs/1206.5716}{arXiv:1206.5716}.
The matrix reduction and its four-row solution above are proved here
and do not require the module-trace theorem.
The full-center argument uses V. Ostrik, \emph{Pivotal fusion categories
of rank 3}, Theorem 2.13, Corollary 2.16 and especially Remark 2.19,
\href{https://arxiv.org/abs/1309.4822}{arXiv:1309.4822}, pp. 6--8.
The center indicator formula is from S.-H. Ng and P. Schauenburg,
\emph{Frobenius--Schur indicators and exponents of spherical categories},
Theorem 4.1, \href{https://arxiv.org/abs/math/0601012}{arXiv:math/0601012}.




